% \subfile{./DocumentazioneRaccolta/Analisi/Raccolta_Requisiti.tex}
\chapter{Raccolta dei requisiti}
\section{Scelta dei candidati}
La raccolta requisiti si è svolta mediante la selezione di figure aziendali con distinti ruoli ed esperienza.
Queste figure selezionate sono state poi sottoposte ad una intervista sia in modalità vis a vis,
che in alcuni casi da remoto per favorire la registrazione del meeting.
I candidati possono poi essere raggruppati in soci aziendali, utilizzatori, e non utilizzatori.
All'interno degli ultimi due gruppi sono state selezionate, per far fronte ai basias dovuti all'abitudine, candidati sia con esperienza pluriennale che
nuovi assunti con poca esperienza limitata.
In totale sono state intervistate 14 persone così suddivise:
\paragraph{Utilizzatori abituali:}
\begin{itemize}
    \item \textbf{Agenti commerciali esterni:} 2 persone
    \begin{itemize}
        \item 1 agente con esperienza limitata.
        \item 1 agente con esperienza consolidata.
    \end{itemize}
    \item \textbf{Commerciali interni:} 4 persone
    \begin{itemize}
        \item 1 specializzato in nuovi impianti.
        \item 1 specializzato in ricambi.
        \item 2 con competenze trasversali sia per quanto concerne nuovi impianti che post-vendita.
    \end{itemize}
\end{itemize}
\paragraph{Utilizzatori saltuari:}
\begin{itemize}
    \item \textbf{Soci aziendali con responsabilità:} 3 persone
    \begin{itemize}
        \item 1 responsabile delle vendite
        \item 1 responsabile tecnico
        \item 1 responsabile amministrativo e CEO
    \end{itemize}
    \item \textbf{Direttore di produzione:} 1 persona
\end{itemize}

\paragraph{Non utilizzatori:}
\begin{itemize}
    \item \textbf{Impiegati tecnici:} 4 persone
    \begin{itemize}
        \item 1 Project manager.
        \item 1 Responsabile logistica.
        \item 1 Progettista meccanico.
        \item 1 Progettista elettrico.
    \end{itemize}
\end{itemize}
L'ordine delle interviste svolte non è stato lineare reparto per reparto, ma si è optato di partire prima di tutto
dalla proprietà per capire la direzione ed aspettative del progetto.
Le restanti interviste hanno seguito un ordine sparso alternando 
utilizzatori e non utilizzatori, cercando di non saturare subito ogni reparto in modo da sfruttare quanto appreso
dalle intervista precedenti per guidare le successive.

\section{Modalità di svolgimento raccolta informazioni}
Le interviste sono state condotte in modalità mista, in presenza e da remoto, secondo la 
disponibilità degli intervistati, ma sempre singolarmente.
Ogni intervista,con permesso degli interlocutori, è stata registrata anche solo l'audio per poi essere rielaborata e 
trascritta anche mediante strumenti IA.
La scelta di utilizzare la modalità fisica o da remoto è stata dettata dalla necessità, per utilizzatori abituali 
di registrare visivamente come utilizzano l'attuale applicativo captando i flussi di utilizzo e 
le principali funzioni attualmente presenti da mantenere.
Si è utilizzato Microsoft Teams per registrare i meeting.
Gli utenti non utilizzatori, o saltuari si è preferito un approccio in presenza
facilitando la comunicazione e l'interazione e, dove possibile, raccogliendo consigli sulle parti 
del porgetto che potevano interessargli.

Le prime interviste sono state supportate da una scaletta di domande, preparate sulla base dell' AS-IS\ref{chap:AsIs} che hanno aiutato
a guidare l'incontro.
Con il progredire delle interviste è stato adottato un approccio più flessibile, portando i risultati
delle interviste precedenti ai nuovi incontri per approfondire aspetti specifici e chiarire dubbi emergenti, ma soprattutto 
distillare quali bisogni erano i più condivisi.
Questo approccio ha consentito di raccogliere quanti più requisiti possibili, ordinandoli per importanza.

\section{Riassunto punti emersi}
In questa sezione elencherò i principali punti emersi dagli incontri che hanno portato ad una analisi dei requisiti.
\subsection{Rispetto all'attuale sistema}
\subsubsection{Funzionalità da mantenere}
Nelle interviste raccolte sono stati sottolineati una serie di funzionalità da mantenere.
Tali funionalità sono:
\begin{itemize}
    \item Gestione delle scontische a cascata. Attualmente si ha la possibilità di inserirle in più punti, ovvero nella singola voce, nel capitolo,
    o dell'offerta stessa.
    \item Gestire all'interno delle voci o capitoli opzionali che non concorrono al prezzo offerto, ma ad uno opzionale.
    \item Gestione di porzioni di contratto raggruppati in capitoli.
    \item I campi del database che attualmente concorrono alla stampa del contratto, sono da mnatenere.
    \item Mantenimento di una reportistica formato PDF del contratto digitale. Da gestire, come ora, il multiligua nel report.
    \item Possibilità di esportare all'\ac{ERP} le voci contrattuali dell'ordine.
\end{itemize}
\subsubsection{Criticità emerse}
Mediante le dimostrazioni di utilizzo si è potuto raccogliere le 
criticità da attenzionare per evitare di ripresentino, come:
\begin{itemize}
    \item L'attaule sistema si blocca, specialemente in condizioni di concorrenza se più utenti utilizzano in contemporanea il file.
    \item Non è presente una gestione centralizzata dei listini ed l'aggiornamento dei codici, attualmente si gestisce un file Excel manualmente.
    \item Mancanza di linee guida nella compilazione delle offerte. Si è avanzata l'ipotesi di un configuratore di prodotto che aiuti l'utente.
    \item Mancanza di gestione dei valori duplicati e controlli sulle chiavi porta a creare clienti e prodotti duplicati.
\end{itemize}
\subsection{Richeste avanzate}
\subsubsection{Gestione ruoli e permessi}
La scelta di coprire quanti più reparti ha fatto emergere la necessità di gestire i ruoli nell'applicativo.
Si è richiesto quindi di attribuire agli utenti dei ruoli che corrispoderanno a livelli di viste e permessi
sulle informazioni contenute nel programma stesso.
I ruoli individuati sono:
\begin{itemize}
    \item Admin con visione sul'intero sistema, con possibiltà esclusiva di modifica ad alcune tabelle fondamentali.
    \item Commerciali interni che possono gestire nell'interezza clienti e contratti.
    \item Commerciali esterni all'azienda, esempio agenti a partita IVA, possono vedere solo i propri clienti e contratti.
    \item Amministrazione contabile, può modificare i dati contabili dei clienti. Questo ruolo può altresi esportare le voci contrattuali all'\ac{ERP} per gestire la fiscalità.
    \item Project Manager degli impianti con possibilità di revisionare una vopia tecnica della vendita, in modo da intergrare il non venduto o internamente trsnettere ai reparti tecnici contratti tecnicamente corretti.
    \item Visualizzatore senza permessi di srittura, ma solo possibilità di consultare i contratti stipulati.
\end{itemize}
\subsubsection{Gestione dei contratti}
Nella creazione dei contratti, core del progetto, si svilupperà il vero salto di qualità.
Ho attenzionato molto questa parte ed ho raccolto delle richieste che possono, anche in fase di primo rilascio,
dare al richiedente un sistema più ottimale.
Tali migliorie sono:
\begin{itemize}
    \item Attaulmente non sono gestisti storici dei dati, si è quindi richiesto una gestione dello storico delle modifiche.
    \item Possibilità di gestire più versioni di uno stesso contratto legati dal medesimo codice progressivo.
    In caso di approvazione serve far decadere le versioni non approvati.
    \item Importare dal \ac{PLM} dati tecnici dei prodotti per eseguire calcoli utili alla contrattazione (Es: Potenza elettrica totale consumata).
    \item Gestione degli stati di vanzamento di un progetto e relative notifiche.
    \item Revisione tecnica del venduto in modo da dare la possibilità di proporre all'occorenza integrazioni contrattuali.
\end{itemize}
